\section{Síntesis y ampliación}

Lo que más vale la pena trasladar de esta clase a otros modelos no es un conjunto de cifras de Mixtral, sino un método de auditoría: primero confirmar el dense baseline y luego describir con claridad el routing algorithm; reportar a la vez total/active/resident; incluir en el costo la communication y el load balance; y, por último, usar routing statistics e intervention para comprobar si los «experts» tienen de verdad una división de trabajo interpretable.

\subsection{Quince conclusiones centrales}

Las siguientes conclusiones se organizan en tres niveles: estructura, sistema e interpretación, para no tomar el nombre del modelo como respuesta.

\begin{enumerate}
  \item Mixtral conserva la attention de Mistral 7B y solo sustituye el MLP de cada capa por top-two-of-eight experts.
  \item GQA reduce los KV heads y la KV cache; SWA limita el attention neighborhood de cada capa.
  \item El Router puntúa cada token de forma independiente, elige el top-$k$, normaliza y combina con pesos los expert outputs.
  \item Cada capa tiene ocho experts; la numeración de los experts es permutable dentro de la capa y no tiene una identidad natural entre capas.
  \item Total parameters, active parameters y resident memory son tres magnitudes distintas.
  \item Los active parameters aproximan sobre todo el per-token compute; no representan la latency completa ni el dollar cost.
  \item El expert parallelism requiere un all-to-all de tokens, y la topology entre nodes puede convertirse en el cuello de botella.
  \item El load imbalance hace que el expert más ocupado determine la tail del step paralelo.
  \item Con lotes grandes y alto tráfico es más fácil agrandar los GEMM de cada expert y diluir el dispatch overhead.
  \item La MLP-knowledge hypothesis es sugerente desde el punto de vista experimental, pero no es un teorema de localización funcional.
  \item Los benchmarks de 2024 muestran que las mayores ganancias están en tareas knowledge-heavy, pero no prueban un nuevo reasoning algorithm.
  \item El discrete routing ofrece una señal que puede contabilizarse, pero la routing frequency no es una explicación semántica ni causal.
  \item Unos domain experts nítidos pueden dejar ociosos a otros experts y reducir el aprovechamiento.
  \item La expert ablation es una intervention importante, pero debe registrar la layer, la gate normalization y el tratamiento del overflow.
  \item MoE, domain tuning y RAG son ejes distintos y combinables; el expert swapping además requiere router adaptation.
\end{enumerate}

\subsection{Una tabla que enlaza estructura, entrenamiento y serving}

Para convertir los conceptos en revisiones de ingeniería, la tabla siguiente indica, para cada etapa, las variables que más conviene observar y los errores de juicio más comunes.

\begin{table}[H]
\centering
\small
\begin{tabularx}{\textwidth}{L{0.16\textwidth} L{0.27\textwidth} X}
\toprule
Etapa & Magnitudes clave observadas & Error de juicio común \\
\midrule
Architecture & $E,k$, shared/expert params, GQA/SWA & Estimar el número de parámetros a partir del nombre del modelo \\
Training & tokens/expert, aux loss, gradient, optimizer state & Calcular la memoria solo con los active params \\
Placement & expert sharding, node topology, HBM & Suponer que los experts no activados no necesitan residir en memoria \\
Serving & lote, all-to-all, imbalance, tail latency & Usar FLOPs en lugar del cost real \\
Interpretation & routing stats, ablation, counterfactual & Ponerle directamente una etiqueta de domain a cada expert \\
Evaluation & task setup, ejes active/total, history date & Tomar el plot de 2024 como ranking actual \\
\bottomrule
\end{tabularx}
\caption{Marco mínimo de auditoría de MoE, desde la arquitectura hasta la interpretación.}
\end{table}

\subsection{MoE release checklist para la práctica}

La siguiente checklist sirve para revisar artículos, model cards o releases internos; su objetivo principal es evitar que solo se muestre el parameter headline.

\begin{lstlisting}[caption={Sparse MoE release checklist}]
report_total_and_active_parameters()
report_expert_count_topk_and_shared_modules()
report_weight_precision_and_resident_memory()
report_tokens_per_expert_and_load_imbalance()
report_expert_parallel_topology_and_all_to_all_cost()
compare_against_dense_models_at_compute_and_memory_budgets()
publish_routing_analysis_with_causal_controls()
document_quantization_offload_and_failure_cases()
\end{lstlisting}

\begin{warningbox}{Al ver «13B active», pregunte de inmediato}
¿Cuántos parámetros totales tiene el checkpoint? ¿Dónde están alojados todos los experts: en una sola tarjeta, en un solo node o repartidos entre nodes? ¿De qué tamaño es el lote? ¿Cuál es la P99 latency? ¿Hay token drop/reroute? Sin esta información, el active parameter headline solo describe una ruta de cómputo parcial.
\end{warningbox}

\subsection{Preguntas de autoevaluación}

Estas preguntas sirven para comprobar si de verdad se dominan las tres cuentas y la routing evidence.

\begin{enumerate}
  \item ¿Por qué Mixtral no es un ensemble de ocho modelos 7B independientes?
  \item ¿Qué recurso reduce GQA y cuál reduce SWA?
  \item Escriba el proceso del top-two routing, desde los logits hasta la weighted sum.
  \item ¿Por qué $8\times 7\mathrm{B}$ no es el valor exacto de los total ni de los active parameters?
  \item ¿Por qué los active parameters no permiten predecir directamente la serving latency?
  \item ¿Qué hace el expert parallel all-to-all en cada una de las dos fases, dispatch y combine?
  \item ¿Por qué el expert más lento controla el tiempo de finalización de un lote?
  \item ¿En qué se diferencian quantization, offload y sparsification?
  \item ¿Qué puede probar un domain routing histogram y qué no?
  \item ¿Por qué unos experts dedicados solo a código pueden reducir el aprovechamiento global?
  \item ¿Qué variables de control debe registrar una expert ablation?
  \item ¿Por qué RAG y MoE no son alternativas entre sí?
  \item ¿Por qué, tras sustituir un expert, todavía hay que entrenar el router?
  \item ¿En qué condiciones más experts podrían, por el contrario, reducir el throughput?
\end{enumerate}

\subsection{Lecturas adicionales}

\begin{itemize}
  \item Jiang et al., \href{https://arxiv.org/abs/2401.04088}{\emph{Mixtral of Experts}}: artículo principal sobre la estructura del modelo, su desempeño y el routing analysis.
  \item Jiang et al., \href{https://arxiv.org/abs/2310.06825}{\emph{Mistral 7B}}: GQA, SWA y el dense baseline.
  \item Shazeer et al., \href{https://arxiv.org/abs/1701.06538}{\emph{Outrageously Large Neural Networks}}: la sparsely-gated MoE layer.
  \item Fedus et al., \href{https://arxiv.org/abs/2101.03961}{\emph{Switch Transformers}}: entrenamiento disperso a gran escala y parallelism.
  \item Al leer estos artículos, registre también total/active parameters, precision, hardware, lote y communication; no se limite a extraer el benchmark headline.
\end{itemize}

\end{document}
